\section{Related work}
\label{sec:related}

\paragraph{Component caching.} The technique is imported from exact model
counting, where a DPLL search that splits the residual formula into
variable-disjoint components and caches their counts is what makes
\#SAT tractable on structured
inputs~\cite{BacchusDalmaoPitassi2003,SangBeameKautz2005}. The transfer is not
mechanical. There the components partition a propositional formula and the
results are multiplied; here they partition a tableau node carrying a symbolic
zone, the results are conjoined, and soundness needs a model-theoretic argument
(Theorem~\ref{thm:decomp}) because the components must be recombined into a
single timed model that agrees with all of them about the passage of time. The
clock-ownership condition has no counterpart in the propositional setting.

\paragraph{What did not transfer.} Clause learning, backjumping and the rest of
the conflict-driven apparatus~\cite{BiereEtAl2009Handbook} require conflicts,
and \S\ref{sec:problem} measures none on the instances that cost the most. We
record this because the reflex when a search is slow is to reach for CDCL, and
here that reflex is measurably wrong.

\paragraph{Zone abstractions.} Coarser abstractions than max-constant
extrapolation are standard, in particular the $LU$
bounds~\cite{BehrmannBLP2006Lower} and their
refinements~\cite{HerbreteauSW2012Better}. They do not apply here as the
construction stands: the split points a clock is tested against are used as
equalities by the delay rule, so a clock's lower and upper bound sets coincide
and $\mathrm{Extra}_{LU}$ collapses to $\mathrm{Extra}_{M}$. Reworking the
splitting to be one-sided is the prerequisite, and we have not done it.
Inclusion-based subsumption of zones, the passed-list test of a reachability
tool, cannot fire either: the several zones a discrete state carries are the
disjoint cells the delay rule produces, never nested ones.

\paragraph{TCTL satisfiability.} The decision problem and its complexity are
classical~\cite{AlurCourcoubetisDill1990,AlurHenzinger1994ReallyTemporal}, with
tree-automata constructions giving matching upper
bounds~\cite{LaTorreNapoli2001TimedTreeAutomata}; discrete-time quantitative
CTL has its own tableau
tradition~\cite{EmersonEtAl1992QuantitativeTemporalReasoning}. The present paper
is not about the decision problem but about the cost of one construction for it,
and everything here applies to any tableau whose nodes carry an independent
set of obligations.
